\section{Test setup and methodology}

Each test was done in the environment that could be best controlled, and it is described here together with
its limits. The figures of the results section come from the logs saved as evidence, and each one can be traced
back to one of those files.

\subsection{Tests on the PC}

\begin{itemize}
  \item \textbf{Fence engine:} \dHostTests{} tests with \dHostChecks{} checks, and a randomized test of
        \dStressCases{} cases with a fixed seed that compares the engine's distance with the haversine distance.
        They are repeated under memory-error and undefined-behavior detectors.
  \item \textbf{Pure C modules} (\texttt{alert\_color}, \texttt{fence\_square}, \texttt{app\_fsm},
        \texttt{demo\_route}, \texttt{bench\_polygon}): each has its own test. Those of \texttt{fence\_square},
        \texttt{demo\_route} and \texttt{bench\_polygon} use the real engine, not a copy. The one of
        \texttt{app\_fsm} covers the 25 combinations of state and event.
  \item \textbf{Cross-compilation:} the engine is compiled for Cortex-M33 and ESP32-S3 with warnings as
        errors. Compiling does not prove that it works on the ESP32-S3.
\end{itemize}

\subsection{Timing of \texttt{geofence\_check()} on the board}

The firmware is built with the benchmark enabled. At boot, before the GNSS and the button start, it builds
regular polygons of 4, 16, 64, 256 and 1024 vertices with a radius of 500~m, and for each one it times a batch
of 40 calls: 20 points that are surely inside and 20 that are surely outside. It reports the total time and
the mean per call, and counts the wrong classifications.

Since the kernel counter runs at \dClockHz~Hz and each reading has a resolution of \dTickUs~$\mu$s, a single
call is not timed but the whole batch. It is repeated in three consecutive boots. The time is that of the
kernel counter, not that of an oscilloscope on a pin.

\subsection{Demo mode}

The board button is pressed and the console is read while the simulated route runs. The route does not use the
real position: it starts from a fictitious center. The log is checked for the state change, the creation of the
fence, the 42 steps with their distance and color, and the return to searching when pressed again.

\subsection{Time to first fix from a cold start}

The modem keeps satellite data in non-volatile memory, so a normal reset is a hot start. To measure a cold
start, the build that deletes that data at every boot is used (only the oscillator offset of the device itself
is kept, as the modem library documentation advises). The board runs on battery, outdoors, without a PC. A
stopwatch measures from power-on until the LED is steady green, which is equivalent to the first fix with a
reported accuracy of 50~m or better. Three boots were done; the third was left on for five more minutes so the
modem would store its data.

It is a manual measurement, with an uncertainty of about 2~s from the reaction of the person holding the
stopwatch, and it has no console log.

\subsection{Position accuracy, static test}

With the normal firmware, which keeps the stored data, the board is left still, flat and face up, and the
console is recorded. The last boot is analyzed. The first minute of fixes is discarded because the receiver is
still converging, and the scatter of the positions around their own mean is computed, after converting them to
meters east and north:

\begin{itemize}
  \item \textbf{CEP50:} the median of the distance to the mean. Half of the fixes fall within that radius.
  \item \textbf{CEP95:} the 95th percentile of that distance.
  \item \textbf{DRMS:} the root of the mean of the squared distances.
  \item \textbf{Reported accuracy:} the estimate the receiver itself gives for each fix (1 sigma, 2D).
\end{itemize}

This measures \emph{precision}, that is, how much the fixes scatter around their mean. It does not measure
absolute error: there was no reference of known coordinates and no RTK receiver, and a constant error that
affects all fixes equally is not seen. The computation is done by a script that checks itself with points of a
known result and writes only derived figures, without coordinates.

\subsection{Conditions of the field tests}

The three cold starts were done in a neighborhood street between houses, with a partly cloudy sky, in the
middle of the morning. The static test was done from a first-floor window facing south, with the window open
and the board practically outside, with a little tree and power cables nearby. Neither was done in an open
field. The height of the obstacles, the percentage of visible sky and the radio interference were not
recorded.
